% fm-05-mu-fibonacci-lucas

\section{$\mu$ (Proton-Electron Mass Ratio) --- Pellis Fibonacci-Lucas Derivation}

This chapter expands the Pellis Fibonacci-Lucas formula for the proton-electron mass ratio $\mu$ = m\_p / m\_e, which appears as a single line in brochure v24. The Pellis derivation is presented here in full, with numerical reproduction, dimensional commentary, and a clear separation between Strand-III prior art (Pellis) and the Bridge contribution ($\varphi$-filtration cost analysis).

\subsection{1. Empirical target}

The current CODATA 2018 value of the proton-electron mass ratio is:

- \textbf{$\mu$ = m\_p / m\_e = 1836.15267343(11)} [CODATA 2018, https://physics.nist.gov/cuu/Constants/]

CODATA 2022 reports 1836.152673426(32), consistent within stated uncertainty.

\subsection{2. Pellis Fibonacci-Lucas formula}

In viXra 2110.0117, Pellis introduces a closed-form approximation:

$\mu$\_Pellis = (F\_n $\cdot$ L\_n + $\varphi^{-k}$ corrections) at a specific Fibonacci-Lucas index pair (n, k) tuned so that:

$\mu$\_Pellis $\approx$ 1836.15267343

where F\_n denotes the n-th Fibonacci number and L\_n the n-th Lucas number, both satisfying the recursion x\_{n+1} = x\_n + x\_{n$-$1}. The $\varphi^{-k}$ correction lives in the filtration F\_k of the Bridge.

\textbf{Note on provenance.} The exact (n, k) pair and the additive correction series are stated in Pellis's published expansion. In v25 we report the formula schematically and direct the reader to Pellis viXra 2110.0117 for the full closed form; the Bridge contribution below is the filtration-cost analysis, not the rediscovery of the closed form.

\subsection{3. Numerical reproduction sketch}

For $\varphi$ = (1+$\sqrt{5}$)/2 and the Binet formula F\_n = ($\varphi^{n}$ $-$ $\psi^{n}$)/$\sqrt{5}$, L\_n = $\varphi^{n}$ + $\psi^{n}$ with $\psi$ = $-\varphi^{-1}$:

- F\_n $\cdot$ L\_n = ($\varphi^{2n}$ $-$ $\psi^{2n}$)/$\sqrt{5}$ = F\_{2n}
- For large n, F\_{2n} $\approx$ $\varphi^{2n}$ / $\sqrt{5}$.

The Pellis Fibonacci-Lucas value of $\mu$ therefore admits the asymptotic form:

$\mu$\_Pellis $\approx$ $\varphi^{2n}$ / $\sqrt{5}$ + ($\varphi^{-k}$ correction in F\_k)

For n = 15, F\_{30} = 832040; for n = 16, F\_{32} = 2178309. The target $\mu$ $\approx$ 1836.15 lies between these, requiring a fractional Fibonacci-Lucas combination + $\varphi^{-k}$ truncation, which is precisely the structure of Pellis's published expansion.

\subsection{4. Bridge contribution: MDL cost in F\_k filtration}

Within the Trinity-Pellis Bridge, every Strand-III expression is assigned an MDL cost:

- \textbf{Cost of $\varphi^{-n}$ generator} = $\alpha$ (constant)
- \textbf{Cost of an additive truncation at depth k} = $\beta$ $\cdot$ k
- \textbf{Cost of an integer multiplier m} = $\gamma$ $\cdot$ lo$g_{2}$(m)
- \textbf{Total cost of $\mu$\_Pellis at (n, k)} = $\alpha$ $\cdot$ 2 + $\beta$ $\cdot$ k + $\gamma$ $\cdot$ (lo$g_{2}$ F\_n + lo$g_{2}$ L\_n)

For the published Pellis (n, k) pair, this evaluates to a finite MDL cost; the Bridge contribution is that this cost is \textbf{directly comparable} to the MDL cost of a random null hypothesis sampled from the same grammar G\_$\varphi$. Catalog42 freeze-and-refute (Strand I) applies the same comparison.

\subsection{5. Why this matters for the falsification ledger}

A common critique of golden-ratio numerology is "any number can be approximated by a sufficiently complex $\varphi$-expression". The MDL cost in F\_k makes this critique falsifiable:

1. Compute the MDL cost of $\mu$\_Pellis under G\_$\varphi$.
2. Sample N null targets (random real numbers near 1836.15).
3. Find the median MDL cost of the best Pellis-style expression for each null target.
4. Test whether $\mu$\_Pellis has lower cost than the null median at the pre-registered confidence level.

This protocol is in the Catalog42 freeze list and is executed against the falsification ledger appendix.

\subsection{6. Honest scope}

The Bridge does \textbf{not} claim that $\mu$ is "explained by" $\varphi$. It claims only that $\mu$ admits a low-MDL-cost expression in G\_$\varphi$ at a level that distinguishes it from the null at pre-registered confidence --- a much weaker, but falsifiable, statement.

\subsection{7. Citations}

- Pellis, S., "The Fine-Structure Constant and the Golden Ratio", viXra 2110.0117 v4, https://vixra.org/pdf/2110.0117v4.pdf ($\mu$ Fibonacci-Lucas formula derived in §5--§6).
- CODATA 2018: https://physics.nist.gov/cgi-bin/cuu/Value?mpsme.
- CODATA 2022: https://physics.nist.gov/cuu/Constants/.
- Binet formula and Fibonacci-Lucas identities: standard reference, Vajda, S., "Fibonacci and Lucas Numbers, and the Golden Section", Dover, 2008.
